\section{Related Work \texorpdfstring{\cred{[1.0 pages]}}{[1.0 pages]}}
\label{sec:related}
\para{Conversational Memory Benchmarks.}
Conversational memory benchmarks have expanded from testing long-transcript fact retrieval to evaluating persistent and adaptive memory behavior.
LoCoMo~\citep{maharana2024evaluating}, LongMemEval~\citep{wu2025longmemeval}, BEAM~\citep{tavakoli2025beam}, PersonaMem-v2~\citep{jiang2025personamemv2}, and MemGallery~\citep{bei2026memgallery} make long-horizon recall, temporal QA, contradictions, and answerability failures concrete over extended dialog histories. EverMemBench~\citep{hu2026evermembench} moves toward richer persistent environments through multi-party collaborative dialog, group-chat histories, evolving decisions, and role-conditioned personas for fine-grained recall, memory awareness, and profile understanding. MemoryBench~\citep{memorybench2025} studies continuous learning from simulated user feedback and behavior logs across multiple domains, languages, and task formats, rather than ego-centric conversational recall over a shared social world.

As mentioned earlier, the above benchmarks are inadequate for evaluating AI agents' capability of working with on-device, personal, and persistent memory, due to the three structural gaps.
Table~\ref{tab:comparison} makes this audit explicit: prior benchmarks cover important pieces of long-term memory evaluation, but \benchmark targets the combined setting of ego-centric, globally-coherent, and activity-dense conversational memory for personal agents.

\begin{table}[t]
\centering
\caption{Comparison of \benchmark with related memory benchmarks along three axes: the structural gaps identified in \S\ref{sec:intro}, corpus statistics, and the six evaluation dimensions. \cmark/\xmark indicate the dimension is/is not evaluated; $\sim$ indicates partial coverage. Footnotes carry per-cell definitions; full per-benchmark derivations are in App.~\ref{app:tab-comparison-details}.}
\label{tab:comparison}
\vspace{4pt}

\scriptsize
\setlength{\tabcolsep}{3pt}
\renewcommand{\arraystretch}{1.2}

\begin{adjustbox}{max width=\textwidth}
\begin{tabular}{@{}l ccccccc@{}}
\toprule
 & \textbf{LoCoMo} & \textbf{LongMemEval} & \textbf{BEAM (10M)} & \textbf{EverMemBench} & \textbf{PersonaMem-v2} & \textbf{MemGallery} & \cellcolor{ourscolor}\textbf{\benchmark} \\
\midrule
\multicolumn{8}{@{}l}{\textit{Structural differentiators (Gaps 1--3).}} \\
\addlinespace[1pt]
Perspective (Gap 1)          & Omniscient & Omniscient & Omniscient & Omniscient & Omniscient & Omniscient & \cellcolor{ourscolor}\textbf{Ego-centric} \\
Global coherence (Gap 2)     & persona-scaffolded pairs & stitched sessions & indep.\ episodes & group chat & persona-driven LLM sim. & synth.\ + clustered sess. & \cellcolor{ourscolor}\textbf{single multi-agent sim} \\
Activity-density (Gap 3)     & 0.3K$^a$   & ---$^b$    & ---$^b$    & ${\sim}$0.07K & ---$^b$    & ---$^b$    & \cellcolor{ourscolor}\textbf{$24.1$K}$^d$ \\
\midrule
\multicolumn{8}{@{}l}{\textit{Statistics.}} \\
\addlinespace[1pt]
Interaction topology         & dyadic & dyadic & dyadic & multi-user & dyadic & dyadic & \cellcolor{ourscolor}\textbf{multi-user} \\
\midrule
\multicolumn{8}{@{}l}{\textit{Evaluation dimensions.}} \\
\addlinespace[1pt]
\DOneDef & \xmark   & \xmark & \xmark   & \xmark   & \xmark & \xmark & \cellcolor{ourscolor}\cmark \\
\DTwoDef & $\sim$ & \xmark & \xmark   & $\sim$ & \xmark & \xmark & \cellcolor{ourscolor}\cmark \\
\DThreeDef (incl.\ temporal sub-types) & \cmark & \cmark & \cmark & \cmark & \cmark & \cmark & \cellcolor{ourscolor}\cmark \\
\DFourDef & \cmark   & \cmark & \cmark   & \cmark   & \cmark & \cmark & \cellcolor{ourscolor}\cmark \\
\DFiveDef (abstain / answer mode split) & $\sim$ & $\sim^g$ & $\sim$ & \xmark & \xmark & $\sim$ & \cellcolor{ourscolor}\cmark \\
\DSixDef & \xmark   & \xmark & \xmark   & \xmark   & \xmark & \xmark & \cellcolor{ourscolor}\cmark \\
\bottomrule
\end{tabular}
\end{adjustbox}

\vspace{4pt}
\raggedright\scriptsize
$^a$LoCoMo does not publish a calendar-day count; the entry reports tokens per speaker per dated session under a session-as-day proxy, derived in App.~\ref{app:tab-comparison-details}.
$^b$``---'' denotes benchmarks that reach million-token scale via session concatenation or independent-episode sampling, with no defined daily interaction rate and (in most cases) no day boundaries.
$^d$The ${\sim}$10K figure is a conservative human-conversation baseline. \benchmark reports the released text-only ego-observed workload: $24.1$K tokens/agent/day on average, excluding metadata, headers, and prompts; this comprises $18.0$K from person--person/social interactions and $6.1$K from person--agent sessions.
$^g$LongMemEval's abstention sub-task covers 30 of the benchmark's 500 total evaluation questions ($6\%$); a never-abstain policy is thus upper-bounded at $94\%$ of the benchmark, so abstention minimally gates headline accuracy, whereas \DFiveID's balanced abstain-mode / answer-mode design forces every system to make an explicit abstain-vs-answer decision on the abstain-mode items.

\vspace{-6pt}
\end{table}


\para{Persistent Memory Systems.}
%%XM \XM{For agentic AI, \st{Memory systems matter for our setting because they turn}personal memory is not a context-window problem, but a backend-design one.} 
Closed-source persistent memory systems such as ChatGPT Memory~\citep{openai2024chatgptmemory} and Claude Memory~\citep{anthropic2025claudememory} show demand for persistent personalization, while open-source or publicly specified stacks such as Mem0~\citep{mem0_2024}, Memobase~\citep{memobase2025}, MemOS~\citep{memos2025}, Zep/Graphiti~\citep{rasmussen2025zep,zep2025graphiti}, and LangMem~\citep{langchain2025langmem} expose memory backends compatible with different readers. 
This fast-growing line of work motivates \benchmark, to enable their evaluation not from a context-window perspective, but treating them as agentic backends dealing with ego-centric, on-device workloads.

\para{Privacy and Permission.}
Recent work around LLM disclosure splits into three threads. \emph{Theory.} Contextual-integrity accounts argue that appropriate disclosure depends on actor, recipient, and norm rather than on the surface fact alone~\citep{nissenbaum2004privacy}, recently re-articulated for agents as a vision for access-control evaluation~\citep{aacvision2025}. \emph{Single-turn / scenario benchmarks} probe whether a model respects social roles and norms on individually-presented disclosure decisions: can-it-keep-a-secret tests~\citep{mireshghallah2024can}, contextual-integrity scenario suites~\citep{privacylens2024,privacybench2024,magpie2025,ariel2025}, and RL post-training on synthetic CI examples~\citep{cireason2025}. \emph{Agentic / pipeline benchmarks} extend this to tool-use and multi-agent traffic, where leakage emerges from cross-tool aggregation, MCP/A2A messages, or inter-agent channels rather than a single response~\citep{agentdam2025,agentscope2026,privacyaction2025,topbench2025,agentleak2026}; on the systems side, two-tier shared/private memory with dynamic access graphs has been proposed as a mechanism but not benchmarked~\citep{collabmemory2025}. None of these works grade an agent on \emph{memory-conditioned} selective disclosure: whether the model first surfaces the right private content from a long-horizon ego-centric memory and then gates it on the current asker's permissions. \benchmark closes this gap (\S\ref{sec:benchmark}) and our results show why: full-local stacks fail this task either by amnesia or by anti-policy disclosure (Findings 1--2, \S\ref{sec:core-findings}).

% ============================================================================
% 3. BENCHMARK DESIGN
